% TOP_PROBLEM: {"id": 3800005, "title": "Point-hyperplane incidences", "category_id": 6, "category": {"id": 6, "name": "geometry", "display_name": "Geometry", "description": "Euclidean and non-Euclidean geometry, geometric structures.", "slug": "geometry", "order_index": 6, "created_at": "2026-07-31T15:26:25.671Z"}, "status": "partially_solved"}
% FIRST_POSTED: null
% ENTRY_KIND: statement_only
% Imported from: batch12.tex; UnsolvedMath ID 3800005
% Compile twice: lualatex 3800005.tex
\documentclass[11pt]{article}
\usepackage{fontspec}
\setmainfont{Latin Modern Roman}
\usepackage{amsmath,amssymb,mathtools,unicode-math}
\setmathfont{Latin Modern Math}
\usepackage{xurl,hyperref}
\hypersetup{hidelinks,pdftitle={UnsolvedMath Problem 3800005}}
\newcommand{\sref}[2]{\hyperlink{source:#1:#2}{[S#2]}}
\newcommand{\cataloguescope}[1]{\paragraph{Mathematical question.}#1}
\begin{document}
% UnsolvedMath ID 3800005; source catalogue code AMR-037-0005
\setcounter{section}{3800004}
\section{Point-hyperplane incidences}\label{3800005}
{\small\sffamily UnsolvedMath ID 3800005 · Geometry}
\subsection{Definitions and mathematical statement}

\paragraph{Shared notation.}$\mathbb N=\{1,2,\ldots\}$, and $\mathbb Z,\mathbb Q,\mathbb R,\mathbb C$ denote the integers, rationals, reals and complex numbers. Any other conventions are specified locally.


Let $d\ge2$ and $n,m,s,t\in\mathbb N$. Let $P$ be a set of $n$ distinct points and $\mathcal H$ a set of $m$ distinct affine hyperplanes in $\mathbb R^d$. Their incidence graph is the bipartite graph with parts $P,\mathcal H$ and an edge $(x,H)$ exactly when $x\in H$. Its incidence count is
\[
 I(P,\mathcal H)=\#\{(x,H)\in P\times\mathcal H:x\in H\}.
\]
The graph is $K_{s,t}$-free if there do not exist $s$ distinct points and $t$ distinct hyperplanes such that all $st$ incidences occur. Here the part of size $s$ is the point part. Put
\[
 F_d(n,m;s,t)=\max\{I(P,\mathcal H):|P|=n,\ |\mathcal H|=m,\ G(P,\mathcal H)\text{ is }K_{s,t}\text{-free}\}.
\]
The maximum is over actual real point--hyperplane configurations; $0\le I\le nm$ makes the range finite.
\paragraph{Statement recorded in the supplied source.}
Determine the maximum $F_d(n,m;s,t)$, in particular its sharp dependence on $n,m$ for fixed $d,s,t$. The distinguished case is
\[
 F_d(n,n;d,d),
\]
with equal numbers of points and hyperplanes and no $d$ points simultaneously incident to $d$ hyperplanes. Bounds for arbitrary bipartite graphs or for configurations over other fields do not by themselves determine this real-geometric maximum. \sref{3800005}{2}
\subsection{Short English statement}
Find the largest number of incidences between real points and hyperplanes when specified complete bipartite subgraphs are forbidden, especially the balanced case with forbidden $K_{d,d}$.
\subsection{Sources}
\begin{description}
\item[\hypertarget{source:3800005:1}{S1}] ULAM.AI and the UnsolvedMath Contributors, UnsolvedMath: A Curated Collection of Open Mathematics Problems, record 3800005 (AMR-037-0005). CC BY 4.0. \url{https://huggingface.co/datasets/ulamai/UnsolvedMath}.
\item[\hypertarget{source:3800005:2}{S2}] Aleksa Milojević, Benny Sudakov and István Tomon, Incidence bounds via extremal graph theory (2024), Introduction and §1.1, pp.1–4. \url{https://arxiv.org/abs/2401.06670}.
\end{description}
\label{end:3800005}
\end{document}
